\documentclass[14pt]{extarticle}

% Configuración (español + tipografía estilo Times)
\usepackage[T1]{fontenc}
\usepackage[utf8]{inputenc}
\usepackage[spanish]{babel}
\usepackage{newtxtext,newtxmath} % estilo Times (similar a Times New Roman)
\usepackage[letterpaper,margin=2.5cm]{geometry}
\usepackage{amsmath}
\usepackage{array}
\usepackage{enumitem}
\setlist[enumerate]{label=(\arabic*), leftmargin=*, itemsep=0.6em}
\setlength{\arraycolsep}{6pt}

\begin{document}
	
	\noindent\textbf{Soluciones detalladas — División y Teoremas de Polinomios}
	
	% =================== DIVISIÓN LARGA ===================
	\section*{División larga de polinomios (4)}
	\begin{enumerate}[label=(\arabic*)]
		\item \(\displaystyle \frac{x^{3}-8}{x-2}\).
		\[
		\begin{aligned}
			&\text{Complete con términos faltantes: }x^{3}+0x^{2}+0x-8.\\
			&x^{3}\div x = x^{2}\ \Rightarrow\ (x-2)\cdot x^{2}=x^{3}-2x^{2}.\\
			&\underline{(x^{3}+0x^{2})-(x^{3}-2x^{2})}=2x^{2}.\\
			&2x^{2}\div x=2x\ \Rightarrow\ (x-2)\cdot 2x=2x^{2}-4x.\\
			&\underline{(2x^{2}+0x)-(2x^{2}-4x)}=4x.\\
			&4x\div x=4\ \Rightarrow\ (x-2)\cdot 4=4x-8.\\
			&\underline{(4x-8)-(4x-8)}=0.
		\end{aligned}
		\]
		\(\boxed{\text{Cociente }=x^{2}+2x+4,\quad \text{Residuo }=0}\).
		\[
		x^{3}-8=(x-2)(x^{2}+2x+4).
		\]
		
		\item \(\displaystyle \frac{2x^{4}-3x^{3}+0x^{2}+5x-6}{x^{2}-x+2}\).
		\[
		\begin{aligned}
			&2x^{4}\div x^{2}=2x^{2}\ \Rightarrow\ (x^{2}-x+2)(2x^{2})=2x^{4}-2x^{3}+4x^{2}.\\
			&\underline{(2x^{4}-3x^{3}+0x^{2})-(2x^{4}-2x^{3}+4x^{2})}=-x^{3}-4x^{2}.\\[2pt]
			&(-x^{3})\div x^{2}=-x\ \Rightarrow\ (x^{2}-x+2)(-x)=-x^{3}+x^{2}-2x.\\
			&\underline{(-x^{3}-4x^{2}+5x)-(-x^{3}+x^{2}-2x)}=-5x^{2}+7x.\\[2pt]
			&(-5x^{2})\div x^{2}=-5\ \Rightarrow\ (x^{2}-x+2)(-5)=-5x^{2}+5x-10.\\
			&\underline{(-5x^{2}+7x-6)-(-5x^{2}+5x-10)}=2x+4.
		\end{aligned}
		\]
		\(\boxed{\text{Cociente }=2x^{2}-x-5,\quad \text{Residuo }=2x+4}\).
		\[
		2x^{4}-3x^{3}+5x-6=(x^{2}-x+2)(2x^{2}-x-5)+(2x+4).
		\]
		
		\item \(\displaystyle \frac{4x^{5}+0x^{4}+x^{2}-7}{\,2x^{2}+1\,}\).
		\[
		\begin{aligned}
			&4x^{5}\div 2x^{2}=2x^{3}\ \Rightarrow\ (2x^{2}+1)(2x^{3})=4x^{5}+2x^{3}.\\
			&\underline{(4x^{5}+0x^{4}+0x^{3})-(4x^{5}+2x^{3})}=-2x^{3}.\\[2pt]
			&(-2x^{3})\div 2x^{2}=-x\ \Rightarrow\ (2x^{2}+1)(-x)=-2x^{3}-x.\\
			&\underline{(-2x^{3}+0x^{2}+0x)-(-2x^{3}-x)}=x.\\[2pt]
			&x\div 2x^{2}=\tfrac{1}{2x}\text{ (no procede en polinomios)}\ \Rightarrow\ \text{bajamos }+x^{2}-7.\\
			&\text{Trabajamos con el término siguiente: }x^{2}\div 2x^{2}=\tfrac12.\\
			&(2x^{2}+1)\!\left(\tfrac12\right)=x^{2}+\tfrac12.\\
			&\underline{(x^{2}+x-7)-(x^{2}+\tfrac12)}=x-\tfrac{15}{2}.
		\end{aligned}
		\]
		\(\boxed{\text{Cociente }=2x^{3}-x+\tfrac12,\quad \text{Residuo }=x-\tfrac{15}{2}}\).
		\[
		4x^{5}+x^{2}-7=(2x^{2}+1)\!\left(2x^{3}-x+\tfrac12\right)+\left(x-\tfrac{15}{2}\right).
		\]
		
		\item \(\displaystyle \frac{-x^{4}+6x^{2}-9x+4}{-x^{2}+3x-2}\).
		\[
		\begin{aligned}
			&(-x^{4})\div(-x^{2})=x^{2}\ \Rightarrow\ (-x^{2}+3x-2)(x^{2})=-x^{4}+3x^{3}-2x^{2}.\\
			&\underline{(-x^{4}+0x^{3}+6x^{2})-(-x^{4}+3x^{3}-2x^{2})}=-3x^{3}+8x^{2}.\\[2pt]
			&(-3x^{3})\div(-x^{2})=3x\ \Rightarrow\ (-x^{2}+3x-2)(3x)=-3x^{3}+9x^{2}-6x.\\
			&\underline{(-3x^{3}+8x^{2}-9x)-(-3x^{3}+9x^{2}-6x)}=-x^{2}-3x.\\[2pt]
			&(-x^{2})\div(-x^{2})=1\ \Rightarrow\ (-x^{2}+3x-2)(1)=-x^{2}+3x-2.\\
			&\underline{(-x^{2}-3x+4)-(-x^{2}+3x-2)}=-6x+6.
		\end{aligned}
		\]
		\(\boxed{\text{Cociente }=x^{2}+3x+1,\quad \text{Residuo }=-6x+6}\).
		\[
		-x^{4}+6x^{2}-9x+4=(-x^{2}+3x-2)(x^{2}+3x+1)+(-6x+6).
		\]
	\end{enumerate}
	
	% =================== DIVISIÓN SINTÉTICA ===================
	\section*{División sintética (4) — divisor lineal \(x-a\)}
	\textit{Convención: fila 1 = coeficientes; fila 2 = productos acumulados; fila 3 = sumas. La última entrada es el residuo.}
	
	\begin{enumerate}[label=(\arabic*), start=5]
		\item \(P(x)=3x^{4}-2x^{3}+0x^{2}+5x-1\), divisor \(x-2\) (\(a=2\)).
		\[
		\begin{array}{r|rrrrr}
			2 & 3 & -2 & 0 & 5 & -1\\
			&   &  6 & 8 & 16 & 42\\ \hline
			& 3 &  4 & 8 & 21 & \boxed{41}
		\end{array}
		\]
		\(\boxed{Q(x)=3x^{3}+4x^{2}+8x+21,\ \ R=41}\), y
		\[
		P(x)=(x-2)\,Q(x)+R.
		\]
		
		\item \(P(x)=x^{5}-4x^{4}+0x^{3}+x^{2}+0x-6\), divisor \(x+1\) (\(a=-1\)).
		\[
		\begin{array}{r|rrrrrr}
			-1 & 1 & -4 & 0 & 1 & 0 & -6\\
			&   & -1 & 5 & -5& 4 & -4\\ \hline
			& 1 & -5 & 5 & -4& 4 & \boxed{-10}
		\end{array}
		\]
		\(\boxed{Q(x)=x^{4}-5x^{3}+5x^{2}-4x+4,\ \ R=-10}\).
		
		\item \(P(x)=2x^{3}-5x^{2}+2x+3\), divisor \(x-3\) (\(a=3\)).
		\[
		\begin{array}{r|rrrr}
			3 & 2 & -5 & 2 & 3\\
			&   &  6 & 3 & 15\\ \hline
			& 2 &  1 & 5 & \boxed{18}
		\end{array}
		\]
		\(\boxed{Q(x)=2x^{2}+x+5,\ \ R=18}\).
		
		\item \(P(x)=4x^{4}+0x^{3}+0x^{2}-5x+2\), divisor \(x-\tfrac{1}{2}\) (\(a=\tfrac12\)).
		\[
		\begin{array}{r|rrrrr}
			\tfrac12 & 4 & 0 & 0 & -5 & 2\\
			&   & 2 & 1 & \tfrac12 & -\tfrac{9}{4}\\ \hline
			& 4 & 2 & 1 & -\tfrac{9}{2} & \boxed{-\tfrac14}
		\end{array}
		\]
		\(\boxed{Q(x)=4x^{3}+2x^{2}+x-\tfrac{9}{2},\ \ R=-\tfrac14}\).
	\end{enumerate}
	
	% =================== TEOREMA DEL RESIDUO ===================
	\section*{Teorema del residuo (2)}
	\begin{enumerate}[label=(\arabic*), start=9]
		\item \(f(x)=2x^{4}-5x^{3}+3x-7\), residuo al dividir entre \(x+3\) (\(a=-3\)):
		\[
		\begin{aligned}
			f(-3)&=2(-3)^{4}-5(-3)^{3}+3(-3)-7\\
			&=2(81)-5(-27)-9-7\\
			&=162+135-16\\
			&=\boxed{281}.
		\end{aligned}
		\]
		
		\item \(f(x)=-x^{5}+4x^{2}-1\), residuo al dividir entre \(x-2\) (\(a=2\)):
		\[
		\begin{aligned}
			f(2)&=-(2)^{5}+4(2)^{2}-1\\
			&=-32+16-1\\
			&=\boxed{-17}.
		\end{aligned}
		\]
	\end{enumerate}
	
	% =================== TEOREMA DEL FACTOR ===================
	\section*{Teorema del factor (2)}
	\begin{enumerate}[label=(\arabic*), start=11]
		\item Verificar si \(x-1\) es factor de \(f(x)=x^{3}-6x^{2}+11x-6\).
		\[
		f(1)=1-6+11-6=0\ \Rightarrow\ \boxed{x-1\ \text{es factor}}.
		\]
		División sintética con \(a=1\):
		\[
		\begin{array}{r|rrrr}
			1 & 1 & -6 & 11 & -6\\
			&   &  1 & -5 & 6\\ \hline
			& 1 & -5 &  6 & 0
		\end{array}
		\]
		\(\Rightarrow\ Q(x)=x^{2}-5x+6=(x-2)(x-3)\).
		\[
		\boxed{f(x)=(x-1)(x-2)(x-3)}.
		\]
		
		\item Verificar si \(x+2\) es factor de \(f(x)=2x^{3}+3x^{2}-8x-12\).
		\[
		f(-2)=2(-8)+3(4)-8(-2)-12=0\ \Rightarrow\ \boxed{x+2\ \text{es factor}}.
		\]
		División sintética con \(a=-2\):
		\[
		\begin{array}{r|rrrr}
			-2 & 2 & 3 & -8 & -12\\
			&   & -4 & 2 & 12\\ \hline
			& 2 & -1 & -6 & 0
		\end{array}
		\]
		\(\Rightarrow\ Q(x)=2x^{2}-x-6=(2x+3)(x-2)\).
		\[
		\boxed{f(x)=(x+2)(x-2)(2x+3)}.
		\]
	\end{enumerate}
	
\end{document}

